% !TEX program = xelatex
\documentclass[12pt,a4paper]{article}

\usepackage[UTF8,heading=true]{ctex}
\setCJKmainfont{SimSun}[BoldFont=SimHei,ItalicFont=KaiTi]
\setCJKsansfont{SimHei}
\usepackage[top=25mm,bottom=25mm,left=25mm,right=25mm,headheight=14pt]{geometry}
\usepackage{amsmath}
\usepackage{enumitem}
\usepackage{xcolor}
\usepackage{hyperref}
\usepackage{fancyhdr}
\hypersetup{colorlinks=true, linkcolor=blue!60!black, urlcolor=blue!60!black}

\pagestyle{fancy}
\fancyhf{}
\fancyhead[L]{\small\color{gray} 《新民主主义论》读书报告}
\fancyhead[R]{\small\color{gray} 姓名/学号待填写}
\fancyfoot[C]{\thepage}

\title{\bfseries 从“中国向何处去”到“建设一个新中国”\\[4pt]\large 《新民主主义论》读书报告}
\author{姓名 / 学号待填写}
\date{2026年4月18日}

\begin{document}
\maketitle

在这次毛概推荐阅读里，我最终选择了《新民主主义论》。最初吸引我的，不只是这篇文章在课程中的重要地位，更是它开篇提出的那个极有压迫感的问题：“中国向何处去”。在1940年的历史条件下，这不是书斋里的抽象追问，而是关系民族生存、革命道路和国家前途的现实难题。重读这篇文章，我最大的感受是，毛泽东并不是先给出一个现成答案，再去套用中国现实，而是先从对中国社会性质的分析入手，一步步把“新民主主义”论证成那个时代最可行、也最有历史方向感的选择。\footnote{毛泽东：《新民主主义论》，《毛泽东选集》第2卷，北京：人民出版社，1991年，第662—706页。}

\section*{一、从时代问题进入理论问题}

我过去在课堂上接触“新民主主义”时，更容易把它理解成几个需要背诵的结论：半殖民地半封建社会、两步走、几个革命阶级联合专政、民族的科学的大众的文化等。但在原文中，这些并不是彼此孤立的知识点，而是一个完整的逻辑链条。毛泽东先指出中国旧政治、旧经济和旧文化之间的关系，进而说明中国革命必须先完成民主主义革命，再向社会主义革命发展。这里最重要的并不是“分两步走”这句判断本身，而是它背后的方法：先认识国情，再决定道路。正因为中国不是一个已经完成资本主义发展的国家，也不是能够直接跨入社会主义的社会，所以革命既不能回到旧式资产阶级民主革命的范畴，也不能脱离实际去空喊“立即社会主义化”。\footnote{毛泽东：《新民主主义论》，《毛泽东选集》第2卷，北京：人民出版社，1991年，第664—667页。}

\section*{二、“新”究竟新在哪里}

我觉得《新民主主义论》里“新”字尤其值得细读。毛泽东强调，新民主主义革命之所以“新”，首先在于它发生在帝国主义时代，已经不再属于旧世界资产阶级民主革命的范畴；其次在于它的领导力量已经发生变化，真正能够把民族独立和社会变革结合起来的是无产阶级及其政党；再次在于它的前途不是资本主义的停顿，而是社会主义的发展方向。正是在这个意义上，新民主主义既不是对旧民主主义的简单重复，也不是对社会主义革命的急躁越级，而是一种立足中国实际、同时又具有世界历史眼光的过渡性安排。读到这里，我会感觉这篇文章的力量，不只是提出了一个新概念，更重要的是重新界定了中国革命在世界历史中的位置。

\section*{三、政治、经济、文化是一整套方案}

这篇文章另一个让我印象很深的地方，是它不是只谈政治口号，而是把政治、经济、文化作为一个整体来设计。原文中提出所谓中华民族的新政治，就是新民主主义的政治；所谓新经济，就是新民主主义的经济；所谓新文化，就是新民主主义的文化。这样的写法让我意识到，毛泽东所设想的“新中国”不是单纯的政权更替，而是一整套社会形态的重建。尤其是“民族的、科学的、大众的文化”这一提法，看起来最像教材中的固定表述，但放在全文语境里，它其实与反帝反封建、动员群众、改造旧文化密切相连。也就是说，新民主主义不是只解决“谁掌权”的问题，它同样关心“怎样建设社会”和“培养什么样的人”。这使得《新民主主义论》具有一种明显的总体性思维，也解释了它为什么能在当时产生广泛影响。

\section*{四、和《中国革命和中国共产党》对读后的理解}

如果把《新民主主义论》同《中国革命和中国共产党》结合起来读，会更清楚地看到它的现实针对性。前者更像是在回答“应该建立一个什么样的新中国”，后者则更系统地分析了中国革命的对象、任务、动力和性质。二者放在一起看，可以发现毛泽东思考问题的一个鲜明特点：不是从抽象教条出发，而是始终围绕中国的社会性质、阶级结构和历史任务展开。对我来说，这一点比结论本身更重要。因为今天回头看，我们未必处在1940年的历史环境中，但“从具体国情出发思考道路问题”这一方法，仍然具有很强的启发性。\footnote{毛泽东：《中国革命和中国共产党》，《毛泽东选集》第2卷，北京：人民出版社，1991年，第636—650页。}

\section*{结语}

当然，读这篇文章时我也会意识到，它首先是一篇处在革命进程中的理论文章，因此它的语言带有鲜明的政治动员色彩，很多判断都服务于当时的历史任务。但正因为如此，它反而更能让人看到理论与现实之间的紧密关系：真正有力量的理论，不是悬在空中的概念，而是在回应时代问题时形成的。总体来看，《新民主主义论》之所以重要，不仅因为它系统阐述了新民主主义革命的基本思想，更因为它把“中国向何处去”的时代追问，转化为一条相对清晰的实践路径。对我而言，这篇文章最大的收获，就是重新理解了“新”不是口号式的新，而是在历史条件、领导力量和发展方向都发生变化之后，对中国革命道路作出的创造性回答。

\end{document}
